\begin{explanationbox}
  \textbf{\textcolor{CExplanation}{为什么会有“左加右减”？}}

  最容易想错的地方是：\textbf{点向左移，横坐标明明要减，为什么函数里反而要加？}
  因为“移动一个点”和“写出新函数”问的不是同一件事。

  \medskip
  \textbf{\textcolor{CExplanation}{第一步：先看得见的变化}}

  以 $y=x^2$ 为例，顶点 $(0,0)$ 向左移 $2$ 格，到了 $(-2,0)$。
  新函数必须在 $x=-2$ 时得到 $y=0$；要让平方等于 $0$，括号里就得是 $0$：
  \[
    x=-2\quad\Longrightarrow\quad x+2=0,
    \qquad \boxed{y=(x+2)^2}.
  \]
  所以\textbf{不是“往左走就加”，而是“新位置要找回原来的输入”。}

  \medskip
  \textbf{\textcolor{CExplanation}{第二步：追踪同一个点，彻底弄懂}}

  原图像上有一点 $P(1,1)$。左移 $2$ 格后，变成 $Q(-1,1)$：横坐标从 $1$ 变为 $-1$，\textbf{高度始终是 $1$}。
  当新图像走到 $x=-1$ 时，要取得原来 $x=1$ 的函数值；而 $-1+2=1$。
  \[
    \underbrace{g(-1)}_{\text{新图像的高度}}
    =\underbrace{f(-1+2)}_{\text{找回原来的横坐标}}
    =f(1)=1.
  \]
  \begin{center}
  \begin{tikzpicture}[x=1.0cm,y=0.78cm,>=stealth,font=\small]
    \draw[->,gray!75] (-3.85,0)--(2.2,0) node[right] {$x$};
    \draw[->,gray!75] (0,-0.3)--(0,2.8) node[above] {$y$};
    \draw[thick,blue!80!black,domain=-1.55:1.55,samples=55] plot (\x,{\x*\x});
    \draw[thick,teal!75!black,domain=-3.55:-0.45,samples=55] plot (\x,{(\x+2)*(\x+2)});
    \draw[densely dashed,gray] (-1,1)--(1,1);
    \fill[blue!80!black] (1,1) circle (2.1pt);
    \fill[teal!75!black] (-1,1) circle (2.1pt);
    \node[right,blue!80!black] at (1.0,1.15) {$P(1,1)$};
    \node[left,teal!75!black] at (-1.05,1.25) {$Q(-1,1)$};
    \draw[->,teal!80!black,very thick] (0.72,1.05)--(-0.72,1.05)
       ;\node at (0,-1.25) {左移 $2$ 格};
    \node[below] at (-1,0) {$-1$};
    \node[below] at (1,0) {$1$};
    \node[left] at (0,.55) {$1$};
    \node[blue!80!black,above right] at (0.65,2.25) {$y=x^2$};
    \node[teal!75!black,above left] at (-2.15,2.25) {$y=(x+2)^2$};
  \end{tikzpicture}

  \end{center}

  \textbf{\textcolor{CExplanation}{第三步：用一句话推广}}

  左移 $a$ 格：新位置是 $x$，原位置就是 $x+a$，所以 $g(x)=f(x+a)$；
  右移 $a$ 格：新位置是 $x$，原位置就是 $x-a$，所以 $g(x)=f(x-a)$。

  \medskip
  \textbf{\textcolor{CExplanation}{什么时候用最快？}}
  \begin{itemize}[leftmargin=*,itemsep=0.35em]
    \item \textbf{给方向、给距离：} 直接把原式中所有 $x$ 整体替换，再按需化简。
    \item \textbf{求平移了多少：} 优先找顶点、对称轴或同一高度的对应点，比展开式子快。
    \item \textbf{有定义域限制：} 图像移到哪，能取的横坐标就跟到哪。
  \end{itemize}

  \textbf{\textcolor{CExplanation}{顿悟一句话：}}点向左走，横坐标减；\textbf{写新函数要找回旧位置，所以 $x$ 加。}
\end{explanationbox}
